\begin{exercise}[Domain of definition]\label{exr:vi38-01}
For a degree-preserving map \(\varphi:S\to T\), prove that \(U_\varphi=\bigcup_{f\in S_+}D_+(\varphi(f))\) is open in \(\Proj T\).
\end{exercise}

\begin{exercise}[Bad contraction]\label{exr:vi38-02}
For \(k[x,y]\to k[u,v]\) with \(x\mapsto u,y\mapsto0\), identify explicitly the points where the Proj map is undefined.
\end{exercise}

\begin{exercise}[Global radical criterion]\label{exr:vi38-03}
Prove that \(T_+\subseteq\sqrt{\varphi(S_+)T}\) implies the induced Proj morphism is defined everywhere.
\end{exercise}

\begin{exercise}[Positive-degree regrading]\label{exr:vi38-04}
Show that a map satisfying \(\varphi(S_r)\subseteq T_{dr}\) becomes degree preserving as \(S\to T^{(d)}\).
\end{exercise}

\begin{exercise}[Base locus of forms]\label{exr:vi38-05}
Show that same-degree forms \(F_0,\ldots,F_m\) define a morphism precisely on the complement of \(V_+(F_0,\ldots,F_m)\).
\end{exercise}

\begin{exercise}[Affine criterion for closed immersion]\label{exr:vi38-06}
Verify directly that a surjection \(A\twoheadrightarrow B\) induces a closed immersion \(\Spec B\hookrightarrow\Spec A\).
\end{exercise}

\begin{exercise}[Quotient chart]\label{exr:vi38-07}
For homogeneous \(I\subset S\), show that \(S_{(f)}\twoheadrightarrow(S/I)_{(\bar f)}\) is surjective with kernel \(I_{(f)}\).
\end{exercise}

\begin{exercise}[Hyperplane inclusion]\label{exr:vi38-08}
Recover the standard inclusion \(\mathbb P^{n-1}\hookrightarrow\mathbb P^n\) from a homogeneous quotient.
\end{exercise}

\begin{exercise}[Closed immersions base change]\label{exr:vi38-09}
Prove affine-locally that closed immersions are stable under arbitrary base change.
\end{exercise}

\begin{exercise}[Projective base change]\label{exr:vi38-10}
Derive base-change stability of projective morphisms from the previous exercise.
\end{exercise}

\begin{exercise}[Quadratic Veronese]\label{exr:vi38-11}
Compute the kernel of \(k[Y_0,Y_1,Y_2]\to k[s,t]^{(2)}\) and recover the conic equation.
\end{exercise}

\begin{exercise}[Veronese generation]\label{exr:vi38-12}
Prove that every monomial of degree \(rd\) is a product of \(r\) monomials of degree \(d\).
\end{exercise}

\begin{exercise}[Veronese dimension]\label{exr:vi38-13}
For \(n=2,d=2\), compute the target dimension \(N\) and list all degree-two coordinates.
\end{exercise}

\begin{exercise}[Segre chart]\label{exr:vi38-14}
Compute the degree-zero ring on \(D_+(z_{ij})\) of the Segre product and identify it with an affine product chart.
\end{exercise}

\begin{exercise}[Segre minors vanish]\label{exr:vi38-15}
Check directly that every \(2\times2\) minor vanishes under \(z_{ij}\mapsto x_i y_j\).
\end{exercise}

\begin{exercise}[Segre kernel]\label{exr:vi38-16}
Complete the monomial-swap proof that the \(2\times2\) minors generate the kernel of the Segre coordinate map.
\end{exercise}

\begin{exercise}[Quadric surface]\label{exr:vi38-17}
Show that the Segre image of \(\mathbb P^1\times\mathbb P^1\) is exactly the quadric \(Z_{00}Z_{11}-Z_{01}Z_{10}=0\).
\end{exercise}

\begin{exercise}[Relative Segre]\label{exr:vi38-18}
Explain why the Segre equations commute with base change from \(\mathbb Z\) to an arbitrary scheme \(Y\).
\end{exercise}

\begin{exercise}[Composition theorem]\label{exr:vi38-19}
Write out the chain of closed immersions used to prove the composition of two projective morphisms is projective.
\end{exercise}

\begin{exercise}[Projectivity of a closed immersion]\label{exr:vi38-20}
Show that every closed immersion \(X\hookrightarrow Y\) is projective using \(\mathbb P^0_Y\cong Y\).
\end{exercise}

\begin{exercise}[Very ample pullback]\label{exr:vi38-21}
For the Veronese embedding \(\nu_d\), prove \(\nu_d^*\OO(1)\cong\OO(d)\) from transition functions.
\end{exercise}

\begin{exercise}[Linear coordinate change]\label{exr:vi38-22}
Show that \(\mathrm{PGL}_{n+1}(k)\) acts on \(\mathbb P^n_k\) through graded automorphisms of the polynomial ring.
\end{exercise}

\begin{exercise}[A rational but not global formula]\label{exr:vi38-23}
Give homogeneous forms with a nonempty base locus and describe the resulting rational map on its maximal domain.
\end{exercise}

\begin{exercise}[Injective points versus closed immersion]\label{exr:vi38-24}
Explain why checking injectivity on \(k\)-points alone cannot prove that a morphism is a closed immersion.
\end{exercise}
